\chapter{Hemorrhoid Banding}

\section*{What Is This Test or Procedure?}
Hemorrhoid banding, or rubber band ligation, is an outpatient procedure that treats internal hemorrhoids by placing a small elastic band around their base to cut off their blood supply.

\section*{Alternative Names}
Rubber band ligation; banding of hemorrhoids; hemorrhoidal ligation

\section*{Principle}
A tight rubber band is applied to the base of an internal hemorrhoid. The banded tissue becomes ischaemic, necroses, and sloughs off within days, reducing the hemorrhoid.

\section*{History}
Rubber band ligation of hemorrhoids was popularized by Barron in the 1960s and remains the most common outpatient treatment for symptomatic internal hemorrhoids.

\section*{Why This Test or Procedure Is Ordered}
Performed for bleeding or prolapsing internal hemorrhoids (grades 1-3) that fail dietary and stool-softening measures.

\section*{Is This a Primary or Secondary Test or Procedure}
Primary treatment for symptomatic internal hemorrhoids.

\section*{How This Test or Procedure Is Performed}
Performed in the office with an anoscope. A ligator applies a band around the hemorrhoid base above the dentate line, where the tissue is pain-insensitive.

\section*{What Preparation Is Needed}
No sedation is required. A bowel preparation is generally unnecessary; the patient's bowel may be emptied.

\section*{Contraindications and Precautions}
External hemorrhoids, bleeding disorders, and anticoagulant use are relative contraindications. Precaution: banding below the dentate line causes severe pain.

\section*{Risks}
Pelvic pain or tenesmus, bleeding, vasovagal reaction, and rarely pelvic infection or septicaemia.

\section*{What Happens After the Test or Procedure}
The patient may feel pressure for a day; mild bleeding is expected when the band sloughs. Stool softeners prevent straining.

\section*{Understanding Your Results}
Improvement is judged by resolution of bleeding and prolapse; several sessions may be needed for multiple hemorrhoids.

\section*{Normal Results}
Symptom relief with no significant bleeding, pain, or infection.

\section*{Follow-up Tests and Procedures}
Follow-up examination; repeat banding for residual hemorrhoids.

\section*{References}
\begin{enumerate}
  \item MedlinePlus. Hemorrhoid banding. U.S. National Library of Medicine; 2025.
  \item Current Medical Diagnosis and Treatment. McGraw-Hill; 2025.
\end{enumerate}

\clearpage
